\section{Siphons of a production model.}\label{s:KB}
\beP[Production zeros determine siphons and first-order terms determine their invasion matrices]
[siphons and transverse invasion][InvF]\label{p:KBsiph}
\label{p:KBresponsezeros}\label{p:KBorder}

Let $\Omega$ be $\mathbb R_+^{n_S}$ or $[0,\bar x]$ with $\bar x\gg0$, and assume
$f$ in \eqref{eq:KB} is $C^1$ on a neighbourhood of $\Omega$ and satisfies
\[
 f_i(x)\ge0\quad(x_i=0,\ x\in\Omega).
\]
In the box case assume also $f_i(x)\le0$ on $x_i=\bar x_i$.
The lower-face inequalities hold automatically when $\Pi\ge0$.
For $\varnothing\ne \Sigma\subseteq\{1,\ldots,n_S\}$ write $z=x_\Sigma$, $y=x_{\Sigma^c}$.
Here $\Sigma$ is a siphon if $\{x\in\Omega:x_\Sigma=0\}$ is forward invariant
(for the lifetime of solutions in the orthant case).
Write $s(M)=\max_{\lambda\in\sigma(M)}\operatorname{Re}\lambda$.
\begin{enumerate}
\item \emph{Face invariance and transverse linearization.}
The set $\Sigma$ is a siphon exactly when
\begin{equation}\label{eq:KB-face-general}
 \Pi_\Sigma(0,y)=V_{\Sigma\Sigma^c}\,y\quad\text{for every }y\text{ with }(0,y)\in\Omega.
\end{equation}
Since $\Pi_\Sigma\ge0$ and $V_{\Sigma\Sigma^c}\le0$, this is equivalent to
\begin{equation}\label{eq:KB-face}
 V_{\Sigma\Sigma^c}=0,\qquad \Pi_\Sigma(0,y)=0\quad\text{for every }y\text{ with }(0,y)\in\Omega.
\end{equation}
On the face,
\[
 J(0,y)=\begin{pmatrix}J^\perp&0\\J_{\mathrm{mix}}&J_{\mathrm{tan}}\end{pmatrix},
 \quad F_\Sigma(y)=\partial_z\Pi_\Sigma(0,y),\quad
 J^\perp=F_\Sigma(y)-V_\Sigma,
\]
where $J_{\mathrm{mix}}=\partial_zf_{\Sigma^c}(0,y)$ and
$J_{\mathrm{tan}}=\partial_yf_{\Sigma^c}(0,y)$.
The matrix $J^\perp$ is Metzler, $F_\Sigma(y)\ge0$, and
\[
 \operatorname{sign}s(J^\perp)
 =\operatorname{sign}\bigl(\rho(F_\Sigma(y)V_\Sigma^{-1})-1\bigr).
\]
At a
boundary equilibrium every eigenvalue of the full Jacobian has negative
real part exactly when this holds for both $J^\perp$ and
$J_{\mathrm{tan}}$. For $\Sigma^c=\varnothing$ the tangential condition is omitted.

\item \emph{Response zeros.} For \eqref{eq:KB-general-response}, suppose
$a_i(0)>0$ for every $i\in \Sigma$ and put
\[
 \mathcal Q_{i,\Sigma}=\{h_i+P_i(0,y):(0,y)\in\Omega\}.
\]
Then the face criterion becomes
\begin{equation}\label{eq:KB-response-zero-test}
 V_{\Sigma\Sigma^c}=0,\qquad
 \mathcal Q_{i,\Sigma}\subseteq\{q:\Phi_i(q)=0\}\quad(i\in \Sigma).
\end{equation}
Keeping $V$, $a_i$ and $P_i$ fixed, replacing responses by
nonnegative functions with the same zero sets
on attained arguments preserves every siphon. Strictly positive responses
on those arguments exclude nonempty siphons when all $a_i(0)>0$.
For an index with $a_i(0)=0$, its production vanishes on the face
regardless of its response.

\item \emph{Nonnegative polynomials $P_i$.} In \eqref{polR},
assume $\Phi_i(q)=0$ for $q\le0$ and $\Phi_i(q)>0$ for $q>0$.
On the full orthant, the criterion is
\[
 V_{\Sigma\Sigma^c}=0,\quad h_i\le0\ (i\in \Sigma),\quad
 \operatorname{supp}\nu\cap \Sigma\ne\varnothing
 \quad(i\in \Sigma,\ a_{i\nu}>0).
\]
On a finite box it is equivalent when $h_i=0$ for $i\in \Sigma$.
For arbitrary $h_i$, the exact box criterion is
\[
 V_{\Sigma\Sigma^c}=0,\qquad h_i+P_i(0,\bar x_{\Sigma^c})\le0\quad(i\in \Sigma).
\]
Here $\operatorname{supp}\nu=\{j:\nu_j>0\}$. For linear $P_i$, zero stimulus,
and $\beta>0$, the monomial condition is $W_{\Sigma\Sigma^c}=0$.
With diagonal $V$, minimal siphons are the source strongly connected
components of the influence graph $j\to i$ when $w_{ij}>0$---the closed
recurrent classes in the opposite row-transition convention for stochastic $W$.

\item \emph{Order in absent variables.}
On a siphon face write
\[
 \Pi_\Sigma(z,y)=\Pi_\Sigma(0,y)+F_\Sigma(y)z+R_\Sigma(z,y).
\]
If $R_\Sigma=O(\|z\|^2)$ at each resident state under consideration,
changing only the remainder preserves $J^\perp$.
More generally, adding a $C^1$ vector field $H$ preserves face dynamics
and the transverse matrix under
\[
 H(0,y)=0,\qquad \partial_zH_\Sigma(0,y)=0.
\]
Preservation of the entire face Jacobian additionally requires
$\partial_zH_{\Sigma^c}(0,y)=0$; otherwise $J_{\mathrm{mix}}$ may change.
These identities must hold along the face, not just at one equilibrium.
In particular, $\Pi_\Sigma(z,y)-\Pi_\Sigma(0,y)=O(\|z\|^2)$ gives
$J^\perp=-V_\Sigma$ along every resident trajectory.
For nonnegative production on a siphon face, $\Pi_\Sigma(0,y)=0$.
\end{enumerate}
\eeP
\begin{proof}
\begin{enumerate}
\item On the face the absent equations are
$\dot z=-V_{\Sigma\Sigma^c}y+\Pi_\Sigma(0,y)$.
Invariance forces this expression to vanish at every initial point
of the face. Conversely, under \eqref{eq:KB-face-general}, a solution
of the resident equations with $z\equiv0$ solves the full system;
uniqueness proves invariance.

Since the off-diagonal entries of $V$ are nonpositive, both
$-V_{\Sigma\Sigma^c}y$ and $\Pi_\Sigma(0,y)$ are nonnegative when production is
nonnegative on the face. Their sum vanishes only if both vanish.
Taking $y=t e_j$ for each resident index and an admissible $t>0$
gives $V_{\Sigma\Sigma^c}=0$. This uses the specified orthant or nondegenerate box.

The triangular form and the Metzler property of $J^\perp$
are Proposition~\ref{p:Tri}. Its one-sided difference quotient
argument, applied to $\Pi_i$ and including $j=i$, proves $F_\Sigma\ge0$ when
production is nonnegative near the face and vanishes on it.
Every principal submatrix $V_\Sigma$ is a nonsingular $M$-matrix;
the nonnegative-production splitting $F_\Sigma-V_\Sigma$ gives the
spectral-radius sign test.
The stability assertion is Proposition~\ref{p:Tri}(2).

\item On the face,
$\Pi_i(0,y)=a_i(0)\Phi_i(h_i+P_i(0,y))$.
For $a_i(0)>0$ its vanishing is precisely the response zero-set
condition; for $a_i(0)=0$ it is automatic. Together with
\eqref{eq:KB-face}, this proves all zero-set assertions.

\item At $y=0$ the argument of $\Phi_i$ equals $h_i$, so $h_i\le0$ is necessary.
Any positive-coefficient monomial containing only resident
variables grows without bound along $y=t\mathbf1$ and eventually
makes this argument positive. Hence such monomials are excluded on
the full orthant. Conversely, every monomial containing an absent
variable vanishes on the face.
On a box the nonnegative polynomial attains its largest resident
value at $\bar x_{\Sigma^c}$, proving the exact box test; for $h_i=0$ every
resident-only monomial is strictly positive there.
For linear $P_i$ and diagonal $V$, the criterion says that no
edge enters $\Sigma$ from outside. Such a set is a union of strongly
connected components and contains a source component of the full
graph. Each source component is itself predecessor closed, so
these and only these are the minimal nonempty siphons.
With nondiagonal $V$, the same proof includes the transfer edges.

\item A differentiable $O(\|z\|^2)$ remainder has zero derivative
at $z=0$. Also $H(0,y)=0$ implies $\partial_yH(0,y)=0$.
These facts identify the preserved entries of the face Jacobian.
\end{enumerate}
\end{proof}

\ssec{Re-splitting the family of the introduction.}
For the family \eqref{eq:intro-family}, write the
bracket as $B_i(x)$ and work on an invariant box $\Omega$ with
$\rho_ix_i\le1$. The production $\Pi_i=(1-\rho_ix_i)B_i\ge0$, but the
nonlinear removal is outside the linear form $-Vx$. Choose
\[
 \mu_i\ge\rho_i\sup_\Omega B_i+\sup_{[0,\bar x_i]}L_i',\qquad
 V'=V+\operatorname{Diag}(\mu_i),\quad \Pi' =\Pi-L+\operatorname{Diag}(\mu_i)x.
\]
Then $\partial_i\Pi_i'\ge0$ and $\Pi_i'(0,x_{-i})=B_i(0,x_{-i})\ge0$,
so $\Pi'\ge0$. The ODE and its upper-face condition are unchanged;
$V'$ remains a nonsingular $M$-matrix. Off-diagonal derivatives are
unchanged, so the Jacobian is Metzler on the box when the responses increase
in the other variables. Capacity factors can make $\partial_i\Pi_i<0$.

For SIS, with $\Pi_i=(1-x_i)H_i$, the
choice $\mu_i\ge\sup H_i$ restores monotonicity on the cube without changing the siphons or the
transverse Jacobians; the value of $R_\Sigma$ depends on the splitting, while its position relative to
one does not. What the siphon calculus uses holds in any case: on a siphon face,
$\Pi_\Sigma(0,y)=0$ and $\Pi\ge0$ give $\partial_z\Pi_\Sigma(0,y)\ge0$ entrywise, the diagonal included.

\beP[Positive networks have invariant boxes and cooperative networks have extremal equilibria]
[existence, bounds, and extremal equilibria][KBpositive]\label{p:KBwell}
Let $\Omega=[0,\bar x]$, $\bar x\gg0$, and suppose $\Pi\ge0$ on $\Omega$ and
$f_i\le0$ on $x_i=\bar x_i$. Then each initial state in $\Omega$ has a unique
global solution remaining in $\Omega$, and an equilibrium exists there.
More generally, if $\Pi$ is defined on the full orthant and $0\le\Pi\le b$,
then without monotonicity all nonnegative solutions are global and
\[
 x(t)\le\bar x+e^{-Vt}(x(0)-\bar x),\qquad
 \bar x=V^{-1}b,\qquad\limsup_{t\to\infty}x(t)\le\bar x.
\]
\begin{enumerate}
\item[(a)] \emph{Cooperative box.} If in addition $J=-V+\partial_x\Pi$ is
Metzler on $\Omega$, the trajectories from $0$ and $\bar x$
converge to the least and greatest equilibria $x^-\le x^+$; every
equilibrium lies between them and every trajectory satisfies
\[
 x^-\le\liminf_{t\to\infty}x(t)\le\limsup_{t\to\infty}x(t)\le x^+.
\]
Consequently a unique equilibrium attracts the whole box, without
requiring contraction.
\item[(b)] \emph{Globally bounded monotone production.} If $0\le\Pi\le b$,
$b\gg0$, and $\partial_x\Pi\ge0$ entrywise (diagonal included) on the
orthant, then each $[0,a\bar x]$, $a\ge1$, is invariant and, for $a>1$,
absorbs every nonnegative solution in finite time. All equilibria lie in
$[0,\bar x]$; the least and greatest are also the limits of
\[
 x^-_0=0,\quad x^+_0=\bar x,\qquad x^\pm_{n+1}=V^{-1}\Pi(x^\pm_n).
\]
\end{enumerate}
\eeP
\begin{proof}
At $x_i=0$,
$f_i=-\sum_{j\ne i}v_{ij}x_j+\Pi_i\ge0$; at $x_i=\bar x_i$ the
assumed inequality is reversed. These inward signs give invariance.
Local uniqueness and compactness then give global solutions.
For existence of an equilibrium, apply Brouwer to the continuous map
$x\mapsto\operatorname{proj}_\Omega(x+\varepsilon f(x))$.
At a fixed point, an interior coordinate forces $f_i=0$;
a zero coordinate forces $f_i\le0$, and an upper coordinate forces
$f_i\ge0$. The inward signs force equality in both boundary cases.

On the orthant, positivity of $e^{-Vt}$ and variation of constants give
\[
 x(t)=e^{-Vt}x(0)+\int_0^t e^{-V(t-s)}\Pi(x(s))\,ds
 \le e^{-Vt}x(0)+V^{-1}(I-e^{-Vt})b.
\]
This bound excludes finite-time blow-up; $e^{-Vt}\to0$ gives
the asserted componentwise upper limit.

In (a), the Metzler Jacobian makes the flow $\varphi_t$ order preserving
\cite{smith1995}, \cite{Hirsch}. Invariance and the semigroup identity give
\[
 \varphi_t(0)\le\varphi_{t+h}(0),\qquad
 \varphi_{t+h}(\bar x)\le\varphi_t(\bar x)\qquad(t,h\ge0).
\]
Both bounded trajectories therefore have limits $x^-,x^+$.
For each $h\ge0$, continuity gives
$\varphi_h(x^\pm)=\lim_{t\to\infty}\varphi_{t+h}(0\text{ or }\bar x)=x^\pm$,
so both limits are equilibria.
Comparison with $\varphi_t(0)$ and $\varphi_t(\bar x)$ bounds every trajectory
and every equilibrium, proving extremality and the limit inequalities.

In (b), $V^{-1}\ge0$ has no zero row, so $\bar x\gg0$.
Since $f(a\bar x)\le-(a-1)b$, cooperativity implies the upper-face
inequalities on $[0,a\bar x]$: if $x_i=a\bar x_i$ and $x\le a\bar x$,
then $f_i(x)\le f_i(a\bar x)\le0$.
For $a>1$, the preceding global bound and decay of $e^{-Vt}$ give
$x(t)\le a\bar x$ for all sufficiently large $t$.
Finally $T=V^{-1}\Pi$ is continuous and increasing, with
$T(0)\ge0$ and $T(\bar x)\le\bar x$. Thus
\[
 0=x_0^-\le x_1^-\le\cdots\le x_n^-\le x_n^+\le\cdots
 \le x_1^+\le x_0^+=\bar x .
\]
Both sequences converge to fixed points of $T$. Every equilibrium
$e$ satisfies $e=T(e)\le\bar x$ and, by induction,
$x_n^-\le e\le x_n^+$; hence the limits are precisely the least
and greatest equilibria.
\end{proof}

\beP[Every positive smooth box model admits the general production representation]
[representation of positive box models][KBrepresentation]
Let $f$ be $C^1$ on a neighbourhood of an invariant box $[0,\bar x]$,
$\bar x\gg0$. Choose
\[
 \mu_i>\max\{0,\sup_{x\in[0,\bar x]}(-\partial_i f_i(x))\},\qquad
 V=\operatorname{Diag}(\mu_i),\qquad \Pi_i(x)=f_i(x)+\mu_i x_i.
\]
Then $\Pi\ge0$ on the box, so the same ODE has a representation with
nonnegative production and inward upper faces. If $f$ is cooperative there, $\partial_x\Pi\ge0$ entrywise,
including the diagonal.
The chosen diagonal coefficients describe a mathematical splitting;
they need not equal the physical removal coefficients of a given model.
\eeP
\begin{proof}
Invariance gives $f_i(0,x_{-i})\ge0$. Integration in the $i$th
coordinate yields
\[
 \Pi_i(x)=f_i(0,x_{-i})+
 \int_0^{x_i}[\partial_i f_i(s,x_{-i})+\mu_i]\,ds\ge0.
\]
For $j\ne i$, $\partial_j\Pi_i=\partial_j f_i$, while
$\partial_i\Pi_i=\partial_i f_i+\mu_i>0$, proving the derivative assertion. The original inward boundary conditions are
unchanged because the vector field itself is unchanged.
\end{proof}
\beR[Smooth ecological and regulatory models]
[ecological models][KBecoG]
Thus the results that assume only nonnegative production on an invariant
box also apply to smooth ecological or regulatory models, even when their natural
production and loss terms have a different form. Global bounds on the
whole orthant and the special threshold bifurcation formulas remain
additional hypotheses.
\eeR

\beR[Smooth onset of a response]
[smooth onset][KBonset]
If $\Phi$ is $C^1$ on all of $\mathbb R$ and vanishes on
$(-\infty,0]$, necessarily $\Phi'(0)=0$. A positive linear onset
$\Phi'(0+)>0$ instead uses a response defined and $C^1$ from the right
on $[0,\infty)$, with a smooth extension for local calculations;
its zero extension to negative arguments is not $C^1$. The results
use smoothness of $\Pi$ on the admissible phase space, in the sense of
Definition~\ref{d:KB}, not an incompatible smoothness assumption on that
zero extension. For the odds responses of Definition~\ref{d:Odd} the
formula $\phi/(1+\phi)$ itself is a smooth extension to a neighbourhood
of $[0,\infty)$, since $1+\phi>0$ near zero.
\eeR

\paragraph{The results using each subclass.}
\begin{enumerate}

\item \emph{Linear $P_i$} are used in the contraction bound \eqref{eq:KB-shared-contraction} (with
$0\le \Phi_i'\le L_i$) and Theorems~\ref{t:KBsharp} (common $\Phi$, $b_i=1$, $h_i=0$, monotone production with a slope bound), \ref{t:KBhom} (common $\Phi$, $V\mathbf1=v\mathbf1$), \ref{t:KBexchange}
(two identical communities, common $\Phi$ and $b$)\xs{ and in
Propositions~\ref{p:EneF} (strictly increasing $\Phi_i$, diagonal
$V$) and~\ref{p:KBdelay} (time-dependent weights, delays)}.

\item \emph{Polynomial $P_i$} are used in Proposition~\ref{p:InvF}
(criterion with $\Phi_i=0$ on $(-\infty,0]$, $\Phi_i>0$ on $(0,\infty)$),
Theorems~\ref{t:siphMulti} ($h_i=0$, common $\Phi$ in the class $\mathcal C_0$)
and~\ref{t:ActF} ($\Phi_i(sP_i)$ with monomials of degree at least two)\nsh{,}
\nsh{coefficient tests for invariant reductions (omitted here)}.
\nsh{Scalar sums $-vs+\sum_kb_k\Phi_k(\alpha_ku+\beta s)$ in
Proposition~\ref{p:KBmixcounter} and Theorem~\ref{t:MixCusp} use linear $P_i$.}

\item \emph{State-dependent factors} \eqref{eq:KB-general-response} are used in
Propositions~\ref{p:InvF} (zero sets of $\Phi_i$ and siphons, with
$a_i(0)>0$)\shortalt{ and}{,} \ref{p:KBshared} (shared responses; $a_i=b_i$ in the contraction
test)\nsh{ and the capacity-factor reduction of the supplementary material (factors constant on groups)}.
\nsh{With $a_i(x_i)=b_i(1-\rho_ix_i)$ and $\Phi_i$ strictly increasing from $\mathbb R$
onto $(0,1)$ it is used in \ifjournal the energy result of the supplementary material\else Proposition~\ref{p:EneF}\fi. }With
$a_i(x_i)=1-x_i$, $\Phi_i(q)=q$, $h_i=0$ and $P_i=H_i$ polynomial it is the SIS
equation\nsh{ \eqref{eq:KB-SIS}, used in
Theorem~\ref{t:KBidentity} and Theorem~\ref{t:Act}, and in the table of SIS results}.
\end{enumerate}

\paragraph{Transverse dynamics and higher-order activation.}

\beR[The transverse matrix of the capacity-factor form and of SIS]
[transverse matrix][KBtransG]
The initial siphon decomposition gives the transverse Metzler property
without full cooperativity; Proposition~\ref{p:InvF} supplies its
production splitting and spectral test. For the
capacity-factor form \eqref{eq:KB-capacity} the factor $1-\rho_ix_i$ equals one
on the face and drops out of the transverse matrix; for SIS contagion,
\[
 J^\perp(y)=A_{\Sigma\Sigma}+\partial_{x_\Sigma}g_\Sigma(0,y)-V_\Sigma.
\]
Only at the origin, or when $\partial_{x_\Sigma}g_\Sigma(0,y)=0$, does this
reduce to $A_{\Sigma\Sigma}-V_\Sigma$. For instance, $g_i=x_jx_k$ with
$i,j\in \Sigma$ and $k\notin \Sigma$ contributes the resident value $y_k$
to the $(i,j)$ transverse entry. Total polynomial degree and degree
in the absent variables must therefore be distinguished.
\eeR

% Preserved original shared-response block; active version moved to commonRed.
% \paragraph{A common response within each siphon is a restriction, not a new ODE.}
% For a specified family $\mathcal S$ of siphons (for example, the minimal
% siphons of a zero-onset model), impose
% \begin{equation}\label{eq:KB-share-siphons}
%  F_i=F_j\quad\text{as functions whenever }i,j\in S,\quad S\in\mathcal S.
% \end{equation}
% Overlaps force agreement. Let $C_1,\ldots,C_m$ be the resulting equivalence
% classes, with uncovered indices initially singleton classes. Equivalently,
% \begin{equation}\label{eq:KB-shared-response}
%  F_i=\mathcal F_{\kappa(i)},\quad i\in C_{\kappa(i)},\qquad
%  \Pi_i(x)=a_i(x_i)\mathcal F_{\kappa(i)}(h_i+P_i(x)).
% \end{equation}
% Covered classes are unions of overlapping siphons and hence siphons:
% their zero faces are intersections of invariant faces. Uncovered classes
% need not be siphons; further response sharing is optional. Sharing on
% \emph{every} siphon would force one global response whenever the full index
% set is a siphon. Preserving the response zero sets preserves the chosen
% family (Proposition~\ref{p:InvF}).
%
% \beP[Sharing responses preserves the general theory but does not imply reduction]
% [shared responses and theorem scope][KBshared]
% Consider \eqref{eq:KB-shared-response} with the positivity and
% invariance assumptions of Proposition~\ref{p:KBpositive}.
% \begin{enumerate}
% \item Every result stated for production models remains valid under
% this restriction when its own hypotheses hold. Sharing does
% not imply cooperativity, irreducibility, contraction or symmetry.
% For constant amplitudes $a_i=b_i>0$, linear inputs
% $P_i=\beta(Wx)_i$, $\beta\ge0$, $W\ge0$, and
% $0\le\mathcal F_a'\le L_a$, a sufficient condition for exponential contraction is
% \begin{equation}\label{eq:KB-shared-contraction}
%  \beta\rho\!\left(
%  \operatorname{Diag}(b_iL_{\kappa(i)})WV^{-1}\right)<1.
% \end{equation}
% \item On a siphon face where $P_i(0,y)=0$ for $i\in S$ and
% $\mathcal F_{\kappa(i)}(h_i)=0$, the transverse matrix is
% \begin{equation}\label{eq:KB-shared-transverse}
%  J^\perp(y)=-V_{SS}+
%  \operatorname{Diag}\!\left(
%  a_i(0)\mathcal F_{\kappa(i)}'(h_i):i\in S\right)
%  \partial_{x_S}P_S(0,y).
% \end{equation}
% Derivatives at a boundary input are understood from the admissible
% side when needed. If the response, stimulus and $a_i(0)$ are common
% on $S$, their scalar factor is common; the remaining matrix need
% not have constant rows or equal eigenvalues.
% \item Sharing responses within $C_a$ gives an exact reduction only
% with additional identities. Let $H_{ia}=1$ for $i\in C_a$ and zero
% otherwise. Sufficient conditions are
% \[
%  VH=H\widehat V,\qquad
%  a_i(r)=\widehat a_a(r),\quad h_i=\widehat h_a,\quad
%  P_i(Hy)=\widehat P_a(y)\quad(i\in C_a).
% \]
% Then variables initially equal within each $C_a$ remain equal,
% and their common values satisfy
% \[
%  \dot y_a=-(\widehat V y)_a+
%  \widehat a_a(y_a)\mathcal F_a(\widehat h_a+\widehat P_a(y)).
% \]
% When the system is cooperative, Theorem~\ref{t:RedF} transfers local
% exponential stability exactly between a reduced equilibrium and its lift.
% Without cooperativity, exact reduction alone gives no such conclusion.
% \end{enumerate}
% \eeP
% \begin{proof}
% For the contraction assertion put
% $B=\beta\operatorname{Diag}(b_iL_{\kappa(i)})W$ and $M=B-V$.
% Here $\beta\ge0$. The spectral inequality gives $s(M)<0$.
% Since $M$ is Metzler, there are $q\gg0$ and $c>0$ such that
% $Mq\le-cq$. Moreover $J(x)\le M$ entrywise and $J(x)$ is Metzler.
% The mean-value formula on the convex box therefore gives, for two solutions,
% \[
%  D^+\|x(t)-\widetilde x(t)\|_{\infty,q}
%  \le-c\|x(t)-\widetilde x(t)\|_{\infty,q},
%  \qquad \|w\|_{\infty,q}=\max_i |w_i|/q_i.
% \]
% Integration proves exponential contraction.
%
% Differentiating the production gives
% \[
%  \partial_j\Pi_i
%  =\delta_{ij}a_i'(x_i)\mathcal F_{\kappa(i)}(h_i+P_i(x))
%  +a_i(x_i)\mathcal F_{\kappa(i)}'(h_i+P_i(x))\partial_jP_i(x).
% \]
% The first term vanishes on the specified face, proving
% \eqref{eq:KB-shared-transverse}.
% For the reduction, direct substitution gives $f(Hy)=H\widehat f(y)$.
% Uniqueness then preserves equality within each group, and
% Theorem~\ref{t:RedF} gives the stability transfer.
% \end{proof}
% \beR[Shared responses and correlated thresholds]
% [shared thresholds][KBsharenote]
% Sharing a response means sharing a threshold \emph{distribution};
% it does not mean that all variables use the same random threshold
% realization or the same noise. Such correlated updates require a
% separate microscopic model. For disjoint minimal siphons of a linear
% network the sharing requirement is particularly simple, but even
% then unequal removal coefficients or unequal row sums can prevent
% reduction. For example, $V=\operatorname{Diag}(1,2)$,
% $P_1=x_2$, $P_2=x_1$, $h=0$, and a common zero-onset response
% have the minimal siphon $\{1,2\}$, whereas at $x_1=x_2=r>0$,
% $f_1=-r+F(r)$ and $f_2=-2r+F(r)$ are unequal.
% \eeR
%
% \beR[Response shape, zero sets and sharpness]
% [shape and zeros][KBshapenote]
% Changing response shapes while preserving their zero sets preserves
% siphons, but can change slopes in \eqref{eq:KB-shared-transverse},
% contraction constants, equilibria, folds and cusps. For a globally $C^1$ zero-onset response,
% Remark~\ref{r:KBonset} gives $\mathcal F_a'(0)=0$; hence zero stimulus
% and inputs vanishing on the face give $J^\perp=-V_{SS}$.
% The common-response sharpness theorem
% \ref{t:KBsharp} already satisfies any prescribed equality among
% responses, but its full-support perturbation removes all nonempty
% siphons. It must not be used to assert arbitrary multistability
% while simultaneously preserving a prescribed nontrivial siphon
% family. For SIS the common identity response is already shared
% by every index; restrictions on its infection coefficients and
% susceptible factors remain essential.
% \eeR
%

